\section{Results}
\label{sec:results}

\subsection{Matched Baselines under Severe Skew}

Every result below is on Dataset~B unless it names Dataset~A.
Every configuration in Fig.~\ref{fig:baseline}(a) shares the corpus, partition draw,
encoders,
and optimizer at $\alpha{=}0.1$, $K{=}5$; all but the native-budget FedMME run also
share the 24-local-epoch budget. Local-only training averages a macro-F1 score of
0.297 across the five clients of one partition. Over
two matched seeds, FedAvg reaches $0.662\!\pm\!0.074$ and FedProx
$0.737\!\pm\!0.085$. Their 0.075 gap sits inside the wider of the two sample
standard deviations, so these runs do not separate the two methods. Our
SCAFFOLD--AdamW adaptation lands at $0.070\!\pm\!0.015$, and the
oscillation in Fig.~\ref{fig:baseline}(b) points at our AdamW control-variate adaptation rather than
at classical SGD-based SCAFFOLD. The FedMME-style baseline trains five client models
independently and combines them by equal-weight hard voting with mean-softmax ties. At our matched 24-epoch budget it reaches
$0.647\!\pm\!0.080$ (seeds 0.703, 0.591), and at FedMME's native 100-epoch budget
it reaches $0.609\!\pm\!0.109$ (seeds 0.687, 0.532). Quadrupling local computation
does not close the gap to iterative aggregation, so the gap belongs to one-shot
aggregation under severe skew, not to a budget handicap. One aggregation has no later round in which to repair a
poor local solution, and both budgets stay below every iterative method. Every configuration except
SCAFFOLD--AdamW clears
local-only training, though the overlapping intervals make that a coarse ordering
rather than a ranking.

\begin{figure}[!t]
\centering
\begin{minipage}[t]{0.425\columnwidth}\centering
\includegraphics[width=\linewidth]{reviewer_matched_results_pa}\\[-5pt]
{\scriptsize\textbf{(a)}}
\end{minipage}\hfill
\begin{minipage}[t]{0.425\columnwidth}\centering
\includegraphics[width=\linewidth]{reviewer_matched_results_pb}\\[-5pt]
{\scriptsize\textbf{(b)}}
\end{minipage}\\[1pt]
\begin{minipage}[t]{0.425\columnwidth}\centering
\includegraphics[width=\linewidth]{legacy_fig5_centralized_vs_federated}\\[-5pt]
{\scriptsize\textbf{(c)}}
\end{minipage}\hfill
\begin{minipage}[t]{0.425\columnwidth}\centering
\includegraphics[width=\linewidth,trim=0 0 875bp 0,clip]{legacy_fig6_federated_convergence}\\[-5pt]
{\scriptsize\textbf{(d)}}
\end{minipage}
\caption{(a) Matched baselines on Dataset~B at $\alpha{=}0.1$, $K{=}5$.
(b) Convergence under severe skew. (c) Centralized against FedAvg on Dataset~A.
(d) Federated convergence on Dataset~A.}
\label{fig:baseline}
\end{figure}

\textit{Dataset~A.} On the synthetic-image corpus, centralized
text/image/concatenation F1-score was 0.934/0.664/0.956 against FedAvg
0.953/0.632/0.956, and multimodal fusion led both unimodal branches. Neither its
text score nor its image score is an independent measurement of a radiograph
task, so we draw no Dataset~B conclusion from it.

\subsection{Anti-Collapse, Fusion, and Initialization Ablations}

Under severe skew the full class-balanced-sampler and entropy-diversity stack
reaches $0.683\!\pm\!0.089$ F1 at minimum predicted-class diversity 0.90. Drop the
sampler, the entropy-diversity term, or both, and F1 rises to
$0.716\!\pm\!0.055$, $0.726\!\pm\!0.035$, and $0.758\!\pm\!0.016$, while minimum
diversity falls to 0.70, 0.80, and 0.70
  (Fig.~\ref{fig:anti_missing}(a)). The stack buys class coverage and costs accuracy, and
its 0.075 F1 deficit to the unregularized configuration sits inside its own seed
spread. The
confidence penalty in \eqref{eq:loss} stays fixed
across all four. At $\alpha{=}1$ the same four configurations land within
$0.809$--$0.846$ with every minimum diversity at 1.0. On real radiographs these
components no longer earn their place on F1. They remain a diversity guard under
severe skew and are inert at moderate skew. One more
observation: three suites nominally repeat the same severe-skew
FedAvg/concatenation setting, yet return $0.643\!\pm\!0.006$, $0.662\!\pm\!0.074$,
and $0.683\!\pm\!0.089$, a 0.040 span that the widest of those spreads covers. We
report them separately rather than pool them or quote the flattering one.
Figure~\ref{fig:diversity} shows the per-epoch class diversity on Dataset~A.

\begin{figure}[!t]
\centering
\includegraphics[width=0.98\columnwidth,height=0.80in,keepaspectratio]{legacy_fig4_class_diversity}
\caption{Per-epoch class diversity, the fraction of the five classes predicted, on
Dataset~A.}
\label{fig:diversity}
\end{figure}

\begin{figure}[!t]
\centering
\begin{minipage}[t]{0.415\columnwidth}\centering
\includegraphics[width=\linewidth]{reviewer_matched_results_pc}\\[-5pt]
{\scriptsize\textbf{(a)}}
\end{minipage}%
\hfill
\begin{minipage}[t]{0.475\columnwidth}\centering
\includegraphics[width=\linewidth]{reviewer_matched_results_pf}\\[-5pt]
{\scriptsize\textbf{(b)}}
\end{minipage}
\caption{(a) Anti-collapse F1 and diversity under severe skew.
(b) P-FIN-style forced-missing-text results.}
\label{fig:anti_missing}
\end{figure}

Figure~\ref{fig:fusion}(a) trains all eight fusion rules
inside the operational federated loop. The
means run from $0.636$ for CLIP-style fusion up to
$0.759$ for the Flamingo-style rule, with per-rule sample standard
deviations of $0.000$--$0.167$. Since the widest interval exceeds the entire
between-rule spread, the sweep cannot rank fusion operators, and we will not. Projected concatenation, our default everywhere else, sits mid-pack at
$0.659\!\pm\!0.000$.

\begin{figure}[!t]
\centering
\begin{minipage}[t]{0.425\columnwidth}\centering
\includegraphics[width=\linewidth]{reviewer_matched_results_pd}\\[-5pt]
{\scriptsize\textbf{(a)}}
\end{minipage}\hfill
\begin{minipage}[t]{0.425\columnwidth}\centering
\includegraphics[width=\linewidth]{reviewer_matched_results_pe}\\[-5pt]
{\scriptsize\textbf{(b)}}
\end{minipage}\\[1pt]
\begin{minipage}[t]{0.425\columnwidth}\centering
\includegraphics[width=\linewidth]{legacy_fig2_fusion_sweep}\\[-5pt]
{\scriptsize\textbf{(c)}}
\end{minipage}\hfill
\begin{minipage}[t]{0.425\columnwidth}\centering
\includegraphics[width=\linewidth]{legacy_fig3_vit_sweep}\\[-5pt]
{\scriptsize\textbf{(d)}}
\end{minipage}
\caption{(a) Fusion rules and (b) initialization controls on Dataset~B.
(c) Fusion sweep and (d) vision-backbone sweep on Dataset~A.}
\label{fig:fusion}
\end{figure}

Figure~\ref{fig:fusion}(b) separates three federated
initializations at $\alpha{=}1$. The all-random configuration
keeps both architectures but initializes every weight from scratch, opening at 0.535
round-1 F1 and closing at $0.647\!\pm\!0.022$. Public
encoders with random task modules, our operational default, open at 0.646 and close
at $0.831\!\pm\!0.001$. Pooled-start reuses one seed-0 checkpoint from a 12-epoch
centralized run for both FL seeds, so its pretraining sits outside the shared
budget; it opens far higher at 0.874 and ends higher at $0.877\!\pm\!0.007$. Public encoders buy roughly 0.18 F1 over random
initialization, and the pooled checkpoint, selected on the reporting split and so a
diagnostic rather than an unbiased estimate, keeps its head start through round~8.
Non-federated abort-off/on controls post $0.899\!\pm\!0.022$ and
$0.871\!\pm\!0.012$ at diversity 1.0, so the branch never fires and the gap is rerun
variation.

Figure~\ref{fig:anti_missing}(b) is our matched P-FIN-style adaptation, where 256-dimensional
$L_2$-normalized features and concatenation plus an MLP replace the source
bidirectional-attention stack. On the primary forced-missing-text metric the four
imputation rules land within 0.018 of each other: zero filling $0.688\!\pm\!0.025$,
deterministic FIN $0.693\!\pm\!0.004$, Gaussian $\beta$-NLL imputation
$0.706\!\pm\!0.010$, and uncertainty-weighted aggregation $0.706\!\pm\!0.089$. The
means order as P-FIN reports, but the uncertainty-weighted rule ties the plain
probabilistic one and is the noisiest of the four, swinging from 0.770 to 0.643
across seeds, so these runs do not establish that modeling uncertainty tightens the
spread.

\subsection{Retrieval-Stage Evaluation for Retrieval-Augmented Generation (RAG)}

Exact inner-product FAISS search~\cite{Johnson2021faiss} over training-fitted
TF--IDF note vectors gives condition-macro averages of 0.540 for top-1 same-label
retrieval accuracy, 0.468 for same-label precision@5, and 0.707 for mean top-1
cosine similarity, with the five condition rows covering all 600 validation
queries (Fig.~\ref{fig:retrieval}(a)). Top-1 accuracy sits well above the 0.2 label prior yet
far below the 0.707 mean similarity, the signature of a template corpus whose
neighbors are lexically close whether or not they share a label. This check
establishes no clinical retrieval quality.

\begin{figure}[!t]
\centering
\begin{minipage}[t]{0.475\columnwidth}\centering
\includegraphics[width=\linewidth]{fig_rag_validated}\\[-5pt]
{\scriptsize\textbf{(a)}}
\end{minipage}\hfill
\begin{minipage}[t]{0.38\columnwidth}\centering
\includegraphics[width=\linewidth]{legacy_fig10_rag_heatmap}\\[-5pt]
{\scriptsize\textbf{(b)}}
\end{minipage}
\caption{(a) Retrieval metrics by condition on Dataset~B.
(b) Top-$k$ class-similarity matrix on Dataset~A.}
\label{fig:retrieval}
\end{figure}

\subsection{Scalability, Communication, and Resource Cost}

Figure~\ref{fig:scale_accuracy}(a) assembles the full Cartesian product of
$K{=}\{3,5,10,20\}$ and $\alpha{=}\{0.1,1,5\}$, and one factor dominates
throughout: label skew, not client count. Every $\alpha{=}5$ cell lies in
$0.849$--$0.916$ and every $\alpha{=}1$ cell in $0.818$--$0.862$, whereas the
$\alpha{=}0.1$ row drops to $0.643$--$0.738$ with far wider spreads, up to
$\pm0.159$ at $K{=}3$. Within a row, moving $K$ by a factor of nearly seven shifts
the mean by at most 0.10 F1, and the $\alpha{=}0.1$ row is not even monotone in
$K$. Our auxiliary $K{=}5$ checks at $\alpha{=}0.3$ and $0.5$ give
$0.814\!\pm\!0.001$ and $0.796\!\pm\!0.037$. Cells report $A/K$, where $A$ counts
clients whose realized shard reaches four samples: allocation is complete except
under severe skew, where $K{=}20$ retains 17 active
clients. Figure~\ref{fig:scale_accuracy}(b) shows the realized per-client class counts behind one
severe-skew cell, where single clients hold entire classes.

\begin{figure}[!t]
\centering
\begin{minipage}[t]{0.425\columnwidth}\centering
\includegraphics[width=\linewidth]{reviewer_systems_scalability_pa}\\[-5pt]
{\scriptsize\textbf{(a)}}
\end{minipage}\hfill
\begin{minipage}[t]{0.425\columnwidth}\centering
\includegraphics[width=\linewidth]{reviewer_systems_scalability_pf}\\[-5pt]
{\scriptsize\textbf{(b)}}
\end{minipage}
\caption{(a) Macro-F1 over client count and skew. (b) One severe-skew
client allocation.}
\label{fig:scale_accuracy}
\end{figure}

Figure~\ref{fig:resources} reports runtime (left) and peak allocated memory (right)
for the timed block on
one shared H100 NVL. Timed sections span 41--66\,s per round and peak memory
4.77--5.92\,GiB. Neither panel is a scaling curve. Sequential
simulation and contention on a shared server, not client count, explain the
spread.

\begin{figure}[!b]
\centering
\begin{minipage}[t]{0.425\columnwidth}\centering
\includegraphics[width=\linewidth]{reviewer_systems_scalability_pb}\\[-5pt]
{\scriptsize\textbf{(a)}}
\end{minipage}\hfill
\begin{minipage}[t]{0.425\columnwidth}\centering
\includegraphics[width=\linewidth]{reviewer_systems_scalability_pc}\\[-5pt]
{\scriptsize\textbf{(b)}}
\end{minipage}
\caption{(a) GPU-synchronized round time. (b) Peak allocated memory.}
\label{fig:resources}
\end{figure}

For $|\theta|$ trainable parameters, $b$ bytes/parameter, $T$ rounds, and
nominal $K$, the plotted bidirectional volume is
\begin{equation}
V_{\mathrm{nom}}=2KT|\theta|b,
\label{eq:comm}
\end{equation}
Here $|\theta|{=}153{,}935{,}621$, $b{=}4$, and $T{=}8$, so a 615,742,484-byte
model state gives nominal bidirectional volumes of 27.5/45.9/91.8/183.5\,GiB at
$K{=}3/5/10/20$ ($1$ GiB$=2^{30}$ bytes). We keep the nominal count when $A{<}K$.
Serialization, secure aggregation~\cite{Ergun2022globecom}, compression, and
algorithm-specific state sit
outside it. Figure~\ref{fig:costs}(a) puts the
operating point plainly. Communication grows linearly in $K$ while accuracy stays
nearly flat in $K$. At $\alpha{=}5$, going from $K{=}3$ to $K{=}20$ costs 0.067 F1
for $6.7\times$ the traffic.

The branch comparison in Fig.~\ref{fig:costs}(b) quantifies the multimodal
overhead on Dataset~B. Text, image, and multimodal F1 are 0.880/0.737/0.906, the
multimodal figure using the residual cross-attention operator our pooled screen
selects; projected concatenation, whose costs the resource rows report, scores
0.813. One-way FP32 model-state sizes are 0.248/0.324/0.573\,GiB, timed round
durations 1.8/4.3/5.6\,min, and peak memory 2.84/4.33/5.92\,GiB. Multimodal costs
$2.3\times$ the model state, $3.1\times$ the wall time and $2.1\times$ the memory of
the text branch, while scoring 0.026 F1 above text and 0.169 above image. Multimodal is the
highest-scoring branch on Dataset~B, at single seed and $\alpha{=}1$.

\begin{figure}[!t]
\centering
\begin{minipage}[t]{0.425\columnwidth}\centering
\includegraphics[width=\linewidth]{reviewer_systems_scalability_pd}\\[-5pt]
{\scriptsize\textbf{(a)}}
\end{minipage}\hfill
\begin{minipage}[t]{0.425\columnwidth}\centering
\includegraphics[width=\linewidth]{reviewer_systems_scalability_pe}\\[-5pt]
{\scriptsize\textbf{(b)}}
\end{minipage}
\caption{(a) Formula-derived communication volume. (b) Modality-branch
cost.}
\label{fig:costs}
\end{figure}

\section{Discussion and Limitations}

Partition sensitivity is the most consistent finding of the study. At
$\alpha{=}0.1$ even identically configured FedAvg reruns vary substantially, with
honest updates spreading by $\pm0.159$ F1. That spread also limits robust
aggregation. Poisoning-resistant aggregation methods are designed to down-weight
anomalous updates~\cite{Barros2022globecom}, but under severe label skew such a rule
may also discard genuine minority-class signal. Skew and robustness to
malicious clients therefore pull against each other, and measuring where the balance
falls would require an attacker model that we did not define here.

Several bounds apply to the rest of the results. The client-count study is a
sequential simulation, so it omits concurrency, latency, stragglers, dropout, and
heterogeneous accelerators. Eq.~\eqref{eq:comm} and Fig.~\ref{fig:costs}(b) report
analytic costs rather than end-to-end traffic. Clinical claims stay bounded by
Section~\ref{sec:corpus}, since our F1 and retrieval scores do not estimate
diagnostic safety, generalization, or deployment readiness, and two seeds make our
spreads descriptive. We leave larger multi-site data, more seeds, native-method
replications, and privacy testing to future work.

\section{Conclusion}

We presented OmniMed-FL, a multimodal federated learning benchmark covering 18
model variants for five-class clinical condition classification. At
$\alpha{=}1$ and $K{=}5$, DistilBERT leads the text encoders at $F1{=}0.934$,
ViT-Base/16 leads the vision encoders at $F1{=}0.664$, and concatenation achieves
the best overall score of $F1{=}0.956$, improving by 0.022 over text and 0.292 over
vision. The anti-collapse stack maintains five-class prediction diversity for the
text and multimodal variants, while federated training retains 99.1\% of
centralized performance; Fed-VLM matches its centralized counterpart,, and Fed-LLM
slightly exceeds it. These findings show that warm-start initialization and
anti-collapse regularization preserve strong multimodal performance in the
evaluated setting. The scalability analysis also shows that label skew matters
more than client count, communication grows linearly with participating clients,
and multimodal fusion incurs additional model-state, runtime, and memory costs.


Future work will use patient-paired multi-institutional radiology data, formal
differential privacy auditing, asynchronous aggregation, and end-to-end
communication measurements. OmniMed-FL provides a practical blueprint for
decentralized multimodal learning, but external clinical validation and safeguards
remain necessary before deployment in telehealth or diagnostic workflows.
